\documentclass[12pt,a4paper]{article}
\usepackage[utf8]{vietnam}
\usepackage{amsmath,amssymb}
\usepackage{tabularx}
\usepackage{tikz}
\usepackage{xcolor}
\usepackage[top=0.45cm, bottom=0.45cm, left=0.6cm, right=0.6cm]{geometry}
\usepackage[dethi]{ex_test}

\setlength{\parskip}{0pt}
\setlength{\parindent}{0pt}

% Dòng kẻ chấm mờ đều đặn chuẩn viết tay (khoảng cách 19.5pt)
\newcommand{\dongke}[1]{%
  \begingroup%
  \setlength{\parskip}{0pt}%
  \baselineskip=19.5pt%
  \vspace{1mm}%
  \foreach \i in {1,...,#1}{%
    \noindent{\color{gray!60}\dotfill}\par%
  }%
  \vspace{0.8mm}%
  \endgroup%
}

% Định nghĩa ô bài tập tự luận trong bảng
\newcommand{\oBaiTap}[3][5]{%
  \vspace{0.8mm}%
  \noindent\textbf{#2} #3\par%
  \vspace{0.8mm}%
  \noindent\textit{Bài giải:}%
  \dongke{#1}%
}

% Câu trắc nghiệm 4 đáp án dàn hàng ngang
\newcommand{\cauTN}[6]{%
  \noindent\textbf{Câu #1.} #2\par\vspace{1mm}%
  \noindent%
  \makebox[0.25\linewidth][l]{\textbf{A.} #3}%
  \makebox[0.25\linewidth][l]{\textbf{B.} #4}%
  \makebox[0.25\linewidth][l]{\textbf{C.} #5}%
  \makebox[0.25\linewidth][l]{\textbf{D.} #6}\par\vspace{1.5mm}%
}

\begin{document}
\pagestyle{empty}

% #############################################################################
% ========================= TOÁN 9 - TUẦN 03 - BUỔI 01: MÃ ĐỀ 101 ============
% ========================= MẶT TRƯỚC (TRANG 1) ===============================
% #############################################################################

\noindent
\begin{tabularx}{\linewidth}{|X|c|}
\hline
\textbf{PHIẾU KIỂM TRA ĐỊNH KỲ TOÁN 9 - TUẦN 03 - BUỔI 01} & \textbf{Mã đề: 101} \\
\textit{Chủ đề: Phương trình quy về phương trình bậc nhất một ẩn} & \textit{Thời gian: 45 phút} \\
\hline
\end{tabularx}

\vspace{1mm}

\noindent
\begin{tabularx}{\linewidth}{|X|}
\hline
\rule[-2.5mm]{0pt}{8mm}\textbf{Họ và tên học sinh:} \dotfill \quad \textbf{Mã HS:} \makebox[3.5cm]{\dotfill} \\
\hline
\end{tabularx}

\vspace{1mm}

\noindent
\begin{tabularx}{\linewidth}{|p{3.2cm}|X|}
\hline
\centering\textbf{ĐIỂM SỐ} & \centering\textbf{LỜI NHẬN XÉT CỦA THẦY CÔ} \tabularnewline
\hline
\rule{0pt}{1.3cm} & \tabularnewline
\hline
\end{tabularx}

\vspace{1.2mm}

\noindent\textbf{BẢNG ĐIỀN ĐÁP ÁN TRẮC NGHIỆM (Mỗi câu đúng 0,5 điểm):}
\vspace{0.8mm}

\noindent
\renewcommand{\arraystretch}{1.35}
\begin{tabularx}{\linewidth}{|c|*{5}{>{\centering\arraybackslash}X|}}
\hline
\textbf{Câu} & \textbf{1} & \textbf{2} & \textbf{3} & \textbf{4} & \textbf{5} \\
\hline
\textbf{Đáp án} & & & & & \\
\hline
\end{tabularx}
\renewcommand{\arraystretch}{1.0}

\vspace{1.5mm}

\noindent\textbf{\large PHẦN I. TRẮC NGHIỆM KHÁCH QUAN (2,5 điểm)}
\vspace{1mm}

\cauTN{1}{Nghiệm của phương trình $(x - 3)(2x + 4) = 0$ là:}{$x = 3$ hoặc $x = -2$}{$x = -3$ hoặc $x = 2$}{$x = 3$ hoặc $x = 2$}{$x = -3$ hoặc $x = -2$}

\cauTN{2}{Điều kiện xác định của phương trình $\dfrac{x + 2}{x - 5} = 0$ là:}{$x \neq -2$}{$x \neq 5$}{$x > 5$}{$x \neq 0$}

\cauTN{3}{Phương trình $(x + 1)^2 - 4 = 0$ có tập nghiệm là:}{$S = \{1; -3\}$}{$S = \{3; -1\}$}{$S = \{1\}$}{$S = \{-3\}$}

\cauTN{4}{Điều kiện xác định của phương trình $\dfrac{2}{x} - \dfrac{3}{x + 1} = 1$ là:}{$x \neq 0$}{$x \neq -1$}{$x \neq 0$ và $x \neq -1$}{$x \neq 1$}

\cauTN{5}{Nghiệm của phương trình $\dfrac{3x - 2}{x + 1} = 2$ là:}{$x = 4$}{$x = 0$}{$x = -4$}{$x = 3$}

\vspace{1mm}

\noindent\textbf{\large PHẦN II. TỰ LUẬN (7,5 điểm)}

\noindent\textbf{Câu 1 (2,5 điểm). Giải các phương trình sau:}
\vspace{0.8mm}

\noindent
\begin{tabularx}{\linewidth}{|X|X|}
\hline
\oBaiTap[15]{a)}{$(2x - 3)(x + 1) + (x + 1)(x - 4) = 0$} & 
\oBaiTap[15]{b)}{$x^2 - 6x + 9 - 4(x - 3) = 0$} \\
\hline
\end{tabularx}

\newpage
% #############################################################################
% ========================= TOÁN 9 - TUẦN 03 - BUỔI 01: MÃ ĐỀ 101 ============
% ========================= MẶT SAU (TRANG 2) =================================
% #############################################################################

\noindent
\begin{tabularx}{\linewidth}{|X|c|}
\hline
\textbf{PHẦN II. TỰ LUẬN (Tiếp theo) - TOÁN 9 TUẦN 03 BUỔI 01} & \textbf{Mã đề: 101} \\
\hline
\end{tabularx}

\vspace{1mm}

\noindent
\begin{tabularx}{\linewidth}{|X|}
\hline
\rule[-2mm]{0pt}{7mm}\textbf{Họ và tên học sinh:} \dotfill \quad \textbf{Mã HS:} \makebox[3.5cm]{\dotfill} \\
\hline
\end{tabularx}

\vspace{1.5mm}

\noindent\textbf{\large Câu 2 (2,5 điểm). Giải các phương trình chứa ẩn ở mẫu sau:}
\vspace{0.8mm}

\noindent
\begin{tabularx}{\linewidth}{|X|X|}
\hline
\oBaiTap[14]{a)}{$\dfrac{2x - 1}{x - 2} = \dfrac{2x + 3}{x + 1}$} & 
\oBaiTap[14]{b)}{$\dfrac{x + 2}{x - 2} - \dfrac{x - 2}{x + 2} = \dfrac{16}{x^2 - 4}$} \\
\hline
\end{tabularx}

\vspace{1.8mm}

\noindent\textbf{\large Câu 3 (2,5 điểm). Bài toán thực tế:}
\vspace{0.8mm}

\noindent
\begin{tabularx}{\linewidth}{|X|}
\hline
\vspace{0.5mm}
Một ô tô dự định đi từ thành phố X đến thành phố Y cách nhau $200\text{ km}$. Sau khi đi được $\dfrac{1}{2}$ quãng đường với vận tốc dự định, ô tô dừng lại nghỉ $15\text{ phút}$. Để đến Y đúng giờ quy định, trên quãng đường còn lại ô tô phải tăng vận tốc thêm $20\text{ km/h}$ so với vận tốc dự định. Tính vận tốc dự định ban đầu của chiếc ô tô đó.\par
\vspace{0.4mm}
\noindent\textit{Bài giải:}\dongke{16}
\\ \hline
\end{tabularx}

\newpage
% #############################################################################
% ========================= TOÁN 9 - TUẦN 03 - BUỔI 01: MÃ ĐỀ 102 ============
% ========================= MẶT TRƯỚC (TRANG 3) ===============================
% #############################################################################

\noindent
\begin{tabularx}{\linewidth}{|X|c|}
\hline
\textbf{PHIẾU KIỂM TRA ĐỊNH KỲ TOÁN 9 - TUẦN 03 - BUỔI 01} & \textbf{Mã đề: 102} \\
\textit{Chủ đề: Phương trình quy về phương trình bậc nhất một ẩn} & \textit{Thời gian: 45 phút} \\
\hline
\end{tabularx}

\vspace{1mm}

\noindent
\begin{tabularx}{\linewidth}{|X|}
\hline
\rule[-2.5mm]{0pt}{8mm}\textbf{Họ và tên học sinh:} \dotfill \quad \textbf{Mã HS:} \makebox[3.5cm]{\dotfill} \\
\hline
\end{tabularx}

\vspace{1mm}

\noindent
\begin{tabularx}{\linewidth}{|p{3.2cm}|X|}
\hline
\centering\textbf{ĐIỂM SỐ} & \centering\textbf{LỜI NHẬN XÉT CỦA THẦY CÔ} \tabularnewline
\hline
\rule{0pt}{1.3cm} & \tabularnewline
\hline
\end{tabularx}

\vspace{1.2mm}

\noindent\textbf{BẢNG ĐIỀN ĐÁP ÁN TRẮC NGHIỆM (Mỗi câu đúng 0,5 điểm):}
\vspace{0.8mm}

\noindent
\renewcommand{\arraystretch}{1.35}
\begin{tabularx}{\linewidth}{|c|*{5}{>{\centering\arraybackslash}X|}}
\hline
\textbf{Câu} & \textbf{1} & \textbf{2} & \textbf{3} & \textbf{4} & \textbf{5} \\
\hline
\textbf{Đáp án} & & & & & \\
\hline
\end{tabularx}
\renewcommand{\arraystretch}{1.0}

\vspace{1.5mm}

\noindent\textbf{\large PHẦN I. TRẮC NGHIỆM KHÁCH QUAN (2,5 điểm)}
\vspace{1mm}

\cauTN{1}{Nghiệm của phương trình $(x + 4)(2x - 6) = 0$ là:}{$x = 4$ hoặc $x = -3$}{$x = -4$ hoặc $x = 3$}{$x = -4$ hoặc $x = -3$}{$x = 4$ hoặc $x = 3$}

\cauTN{2}{Điều kiện xác định của phương trình $\dfrac{x - 3}{x + 4} = 0$ là:}{$x \neq 3$}{$x \neq -4$}{$x > -4$}{$x \neq 0$}

\cauTN{3}{Phương trình $(x - 2)^2 - 9 = 0$ có tập nghiệm là:}{$S = \{5; -1\}$}{$S = \{1; -5\}$}{$S = \{5\}$}{$S = \{-1\}$}

\cauTN{4}{Điều kiện xác định của phương trình $\dfrac{3}{x} + \dfrac{1}{x - 2} = 2$ là:}{$x \neq 0$ và $x \neq 2$}{$x \neq 0$}{$x \neq 2$}{$x \neq -2$}

\cauTN{5}{Nghiệm của phương trình $\dfrac{4x - 1}{x + 2} = 3$ là:}{$x = 7$}{$x = -7$}{$x = 5$}{$x = 1$}

\vspace{1mm}

\noindent\textbf{\large PHẦN II. TỰ LUẬN (7,5 điểm)}

\noindent\textbf{Câu 1 (2,5 điểm). Giải các phương trình sau:}
\vspace{0.8mm}

\noindent
\begin{tabularx}{\linewidth}{|X|X|}
\hline
\oBaiTap[15]{a)}{$(3x - 2)(x + 2) + (x + 2)(x - 5) = 0$} & 
\oBaiTap[15]{b)}{$x^2 - 8x + 16 - 3(x - 4) = 0$} \\
\hline
\end{tabularx}

\newpage
% #############################################################################
% ========================= TOÁN 9 - TUẦN 03 - BUỔI 01: MÃ ĐỀ 102 ============
% ========================= MẶT SAU (TRANG 4) =================================
% #############################################################################

\noindent
\begin{tabularx}{\linewidth}{|X|c|}
\hline
\textbf{PHẦN II. TỰ LUẬN (Tiếp theo) - TOÁN 9 TUẦN 03 BUỔI 01} & \textbf{Mã đề: 102} \\
\hline
\end{tabularx}

\vspace{1mm}

\noindent
\begin{tabularx}{\linewidth}{|X|}
\hline
\rule[-2mm]{0pt}{7mm}\textbf{Họ và tên học sinh:} \dotfill \quad \textbf{Mã HS:} \makebox[3.5cm]{\dotfill} \\
\hline
\end{tabularx}

\vspace{1.5mm}

\noindent\textbf{\large Câu 2 (2,5 điểm). Giải các phương trình chứa ẩn ở mẫu sau:}
\vspace{0.8mm}

\noindent
\begin{tabularx}{\linewidth}{|X|X|}
\hline
\oBaiTap[14]{a)}{$\dfrac{2x + 1}{x - 1} = \dfrac{2x - 3}{x - 3}$} & 
\oBaiTap[14]{b)}{$\dfrac{x + 3}{x - 3} - \dfrac{x - 3}{x + 3} = \dfrac{24}{x^2 - 9}$} \\
\hline
\end{tabularx}

\vspace{1.8mm}

\noindent\textbf{\large Câu 3 (2,5 điểm). Bài toán thực tế:}
\vspace{0.8mm}

\noindent
\begin{tabularx}{\linewidth}{|X|}
\hline
\vspace{0.5mm}
Một ô tô dự định đi từ thành phố X đến thành phố Y cách nhau $200\text{ km}$. Sau khi đi được $\dfrac{1}{2}$ quãng đường với vận tốc dự định, ô tô dừng lại nghỉ $15\text{ phút}$. Để đến Y đúng giờ quy định, trên quãng đường còn lại ô tô phải tăng vận tốc thêm $20\text{ km/h}$ so với vận tốc dự định. Tính vận tốc dự định ban đầu của chiếc ô tô đó.\par
\vspace{0.4mm}
\noindent\textit{Bài giải:}\dongke{16}
\\ \hline
\end{tabularx}

\newpage
% #############################################################################
% ========================= TOÁN 9 - TUẦN 03 - BUỔI 01: MÃ ĐỀ 103 ============
% ========================= MẶT TRƯỚC (TRANG 5) ===============================
% #############################################################################

\noindent
\begin{tabularx}{\linewidth}{|X|c|}
\hline
\textbf{PHIẾU KIỂM TRA ĐỊNH KỲ TOÁN 9 - TUẦN 03 - BUỔI 01} & \textbf{Mã đề: 103} \\
\textit{Chủ đề: Phương trình quy về phương trình bậc nhất một ẩn} & \textit{Thời gian: 45 phút} \\
\hline
\end{tabularx}

\vspace{1mm}

\noindent
\begin{tabularx}{\linewidth}{|X|}
\hline
\rule[-2.5mm]{0pt}{8mm}\textbf{Họ và tên học sinh:} \dotfill \quad \textbf{Mã HS:} \makebox[3.5cm]{\dotfill} \\
\hline
\end{tabularx}

\vspace{1mm}

\noindent
\begin{tabularx}{\linewidth}{|p{3.2cm}|X|}
\hline
\centering\textbf{ĐIỂM SỐ} & \centering\textbf{LỜI NHẬN XÉT CỦA THẦY CÔ} \tabularnewline
\hline
\rule{0pt}{1.3cm} & \tabularnewline
\hline
\end{tabularx}

\vspace{1.2mm}

\noindent\textbf{BẢNG ĐIỀN ĐÁP ÁN TRẮC NGHIỆM (Mỗi câu đúng 0,5 điểm):}
\vspace{0.8mm}

\noindent
\renewcommand{\arraystretch}{1.35}
\begin{tabularx}{\linewidth}{|c|*{5}{>{\centering\arraybackslash}X|}}
\hline
\textbf{Câu} & \textbf{1} & \textbf{2} & \textbf{3} & \textbf{4} & \textbf{5} \\
\hline
\textbf{Đáp án} & & & & & \\
\hline
\end{tabularx}
\renewcommand{\arraystretch}{1.0}

\vspace{1.5mm}

\noindent\textbf{\large PHẦN I. TRẮC NGHIỆM KHÁCH QUAN (2,5 điểm)}
\vspace{1mm}

\cauTN{1}{Nghiệm của phương trình $(2x - 8)(x + 5) = 0$ là:}{$x = 4$ hoặc $x = -5$}{$x = -4$ hoặc $x = 5$}{$x = 4$ hoặc $x = 5$}{$x = -4$ hoặc $x = -5$}

\cauTN{2}{Điều kiện xác định của phương trình $\dfrac{2x + 3}{x - 4} = 0$ là:}{$x \neq -\dfrac{3}{2}$}{$x \neq 4$}{$x > 4$}{$x \neq 0$}

\cauTN{3}{Phương trình $(x + 3)^2 - 16 = 0$ có tập nghiệm là:}{$S = \{1; -7\}$}{$S = \{7; -1\}$}{$S = \{1\}$}{$S = \{-7\}$}

\cauTN{4}{Điều kiện xác định của phương trình $\dfrac{1}{x} + \dfrac{2}{x - 3} = 4$ là:}{$x \neq 0$ và $x \neq 3$}{$x \neq 0$}{$x \neq 3$}{$x \neq -3$}

\cauTN{5}{Nghiệm của phương trình $\dfrac{5x - 3}{x + 1} = 3$ là:}{$x = 3$}{$x = -3$}{$x = 2$}{$x = 0$}

\vspace{1mm}

\noindent\textbf{\large PHẦN II. TỰ LUẬN (7,5 điểm)}

\noindent\textbf{Câu 1 (2,5 điểm). Giải các phương trình sau:}
\vspace{0.8mm}

\noindent
\begin{tabularx}{\linewidth}{|X|X|}
\hline
\oBaiTap[15]{a)}{$(2x - 5)(x + 3) + (x + 3)(x - 1) = 0$} & 
\oBaiTap[15]{b)}{$x^2 - 10x + 25 - 2(x - 5) = 0$} \\
\hline
\end{tabularx}

\newpage
% #############################################################################
% ========================= TOÁN 9 - TUẦN 03 - BUỔI 01: MÃ ĐỀ 103 ============
% ========================= MẶT SAU (TRANG 6) =================================
% #############################################################################

\noindent
\begin{tabularx}{\linewidth}{|X|c|}
\hline
\textbf{PHẦN II. TỰ LUẬN (Tiếp theo) - TOÁN 9 TUẦN 03 BUỔI 01} & \textbf{Mã đề: 103} \\
\hline
\end{tabularx}

\vspace{1mm}

\noindent
\begin{tabularx}{\linewidth}{|X|}
\hline
\rule[-2mm]{0pt}{7mm}\textbf{Họ và tên học sinh:} \dotfill \quad \textbf{Mã HS:} \makebox[3.5cm]{\dotfill} \\
\hline
\end{tabularx}

\vspace{1.5mm}

\noindent\textbf{\large Câu 2 (2,5 điểm). Giải các phương trình chứa ẩn ở mẫu sau:}
\vspace{0.8mm}

\noindent
\begin{tabularx}{\linewidth}{|X|X|}
\hline
\oBaiTap[14]{a)}{$\dfrac{2x - 3}{x - 1} = \dfrac{2x + 1}{x + 2}$} & 
\oBaiTap[14]{b)}{$\dfrac{x + 1}{x - 1} - \dfrac{x - 1}{x + 1} = \dfrac{8}{x^2 - 1}$} \\
\hline
\end{tabularx}

\vspace{1.8mm}

\noindent\textbf{\large Câu 3 (2,5 điểm). Bài toán thực tế:}
\vspace{0.8mm}

\noindent
\begin{tabularx}{\linewidth}{|X|}
\hline
\vspace{0.5mm}
Một ô tô dự định đi từ thành phố X đến thành phố Y cách nhau $200\text{ km}$. Sau khi đi được $\dfrac{1}{2}$ quãng đường với vận tốc dự định, ô tô dừng lại nghỉ $15\text{ phút}$. Để đến Y đúng giờ quy định, trên quãng đường còn lại ô tô phải tăng vận tốc thêm $20\text{ km/h}$ so với vận tốc dự định. Tính vận tốc dự định ban đầu của chiếc ô tô đó.\par
\vspace{0.4mm}
\noindent\textit{Bài giải:}\dongke{16}
\\ \hline
\end{tabularx}

\newpage
% #############################################################################
% ========================= TOÁN 9 - TUẦN 03 - BUỔI 01: MÃ ĐỀ 104 ============
% ========================= MẶT TRƯỚC (TRANG 7) ===============================
% #############################################################################

\noindent
\begin{tabularx}{\linewidth}{|X|c|}
\hline
\textbf{PHIẾU KIỂM TRA ĐỊNH KỲ TOÁN 9 - TUẦN 03 - BUỔI 01} & \textbf{Mã đề: 104} \\
\textit{Chủ đề: Phương trình quy về phương trình bậc nhất một ẩn} & \textit{Thời gian: 45 phút} \\
\hline
\end{tabularx}

\vspace{1mm}

\noindent
\begin{tabularx}{\linewidth}{|X|}
\hline
\rule[-2.5mm]{0pt}{8mm}\textbf{Họ và tên học sinh:} \dotfill \quad \textbf{Mã HS:} \makebox[3.5cm]{\dotfill} \\
\hline
\end{tabularx}

\vspace{1mm}

\noindent
\begin{tabularx}{\linewidth}{|p{3.2cm}|X|}
\hline
\centering\textbf{ĐIỂM SỐ} & \centering\textbf{LỜI NHẬN XÉT CỦA THẦY CÔ} \tabularnewline
\hline
\rule{0pt}{1.3cm} & \tabularnewline
\hline
\end{tabularx}

\vspace{1.2mm}

\noindent\textbf{BẢNG ĐIỀN ĐÁP ÁN TRẮC NGHIỆM (Mỗi câu đúng 0,5 điểm):}
\vspace{0.8mm}

\noindent
\renewcommand{\arraystretch}{1.35}
\begin{tabularx}{\linewidth}{|c|*{5}{>{\centering\arraybackslash}X|}}
\hline
\textbf{Câu} & \textbf{1} & \textbf{2} & \textbf{3} & \textbf{4} & \textbf{5} \\
\hline
\textbf{Đáp án} & & & & & \\
\hline
\end{tabularx}
\renewcommand{\arraystretch}{1.0}

\vspace{1.5mm}

\noindent\textbf{\large PHẦN I. TRẮC NGHIỆM KHÁCH QUAN (2,5 điểm)}
\vspace{1mm}

\cauTN{1}{Nghiệm của phương trình $(3x - 9)(x + 2) = 0$ là:}{$x = 3$ hoặc $x = -2$}{$x = -3$ hoặc $x = 2$}{$x = 3$ hoặc $x = 2$}{$x = -3$ hoặc $x = -2$}

\cauTN{2}{Điều kiện xác định của phương trình $\dfrac{3x - 1}{x + 6} = 0$ là:}{$x \neq -6$}{$x \neq \dfrac{1}{3}$}{$x > -6$}{$x \neq 0$}

\cauTN{3}{Phương trình $(x - 1)^2 - 25 = 0$ có tập nghiệm là:}{$S = \{6; -4\}$}{$S = \{-6; 4\}$}{$S = \{6\}$}{$S = \{-4\}$}

\cauTN{4}{Điều kiện xác định của phương trình $\dfrac{5}{x} - \dfrac{2}{x - 4} = 3$ là:}{$x \neq 0$ và $x \neq 4$}{$x \neq 0$}{$x \neq 4$}{$x \neq -4$}

\cauTN{5}{Nghiệm của phương trình $\dfrac{6x - 5}{x + 1} = 4$ là:}{$x = \dfrac{9}{2}$}{$x = -\dfrac{9}{2}$}{$x = 3$}{$x = 1$}

\vspace{1mm}

\noindent\textbf{\large PHẦN II. TỰ LUẬN (7,5 điểm)}

\noindent\textbf{Câu 1 (2,5 điểm). Giải các phương trình sau:}
\vspace{0.8mm}

\noindent
\begin{tabularx}{\linewidth}{|X|X|}
\hline
\oBaiTap[15]{a)}{$(3x - 4)(x + 1) + (x + 1)(x - 2) = 0$} & 
\oBaiTap[15]{b)}{$x^2 - 12x + 36 - 5(x - 6) = 0$} \\
\hline
\end{tabularx}

\newpage
% #############################################################################
% ========================= TOÁN 9 - TUẦN 03 - BUỔI 01: MÃ ĐỀ 104 ============
% ========================= MẶT SAU (TRANG 8) =================================
% #############################################################################

\noindent
\begin{tabularx}{\linewidth}{|X|c|}
\hline
\textbf{PHẦN II. TỰ LUẬN (Tiếp theo) - TOÁN 9 TUẦN 03 BUỔI 01} & \textbf{Mã đề: 104} \\
\hline
\end{tabularx}

\vspace{1mm}

\noindent
\begin{tabularx}{\linewidth}{|X|}
\hline
\rule[-2mm]{0pt}{7mm}\textbf{Họ và tên học sinh:} \dotfill \quad \textbf{Mã HS:} \makebox[3.5cm]{\dotfill} \\
\hline
\end{tabularx}

\vspace{1.5mm}

\noindent\textbf{\large Câu 2 (2,5 điểm). Giải các phương trình chứa ẩn ở mẫu sau:}
\vspace{0.8mm}

\noindent
\begin{tabularx}{\linewidth}{|X|X|}
\hline
\oBaiTap[14]{a)}{$\dfrac{2x - 1}{x + 1} = \dfrac{2x + 5}{x + 5}$} & 
\oBaiTap[14]{b)}{$\dfrac{x + 2}{x - 2} - \dfrac{x - 2}{x + 2} = \dfrac{8}{x^2 - 4}$} \\
\hline
\end{tabularx}

\vspace{1.8mm}

\noindent\textbf{\large Câu 3 (2,5 điểm). Bài toán thực tế:}
\vspace{0.8mm}

\noindent
\begin{tabularx}{\linewidth}{|X|}
\hline
\vspace{0.5mm}
Một ô tô dự định đi từ thành phố X đến thành phố Y cách nhau $200\text{ km}$. Sau khi đi được $\dfrac{1}{2}$ quãng đường với vận tốc dự định, ô tô dừng lại nghỉ $15\text{ phút}$. Để đến Y đúng giờ quy định, trên quãng đường còn lại ô tô phải tăng vận tốc thêm $20\text{ km/h}$ so với vận tốc dự định. Tính vận tốc dự định ban đầu của chiếc ô tô đó.\par
\vspace{0.4mm}
\noindent\textit{Bài giải:}\dongke{16}
\\ \hline
\end{tabularx}

\end{document}
